% Source-aligned English for printed p.141.
\SourcePage{141}

\LayerTarget{tohoku:remark:satellites-existence:p2}
\noindent that every covariant additive functor $F:C\longrightarrow C'$ has a
universal $\partial$-functor defined in all degrees and equal to $F$ in degree
zero, that is, that all satellites $S^iF$ exist, it plainly suffices to prove
that $S^1F$ and $S^{-1}F$ exist, by the formulas
\[
 S^1(S^iF)=S^{i+1}F\quad\hbox{for }i\geq0,\qquad
 S^{-1}(S^iF)=S^{i-1}F\quad\hbox{for }i\leq0,
\]
\noindent whose verification directly from the definition is trivial. The
problems of finding $S^1F$ and $S^{-1}F$ are dual: these two functors are
interchanged when $C$ and $C'$ are replaced by their opposite categories.

\LayerTarget{tohoku:def:effaceable-coeffaceable}
An additive functor $F:C\longrightarrow C'$ is called \textit{effaceable} if,
for every $A\in C$, there exists a monomorphism $u:A\longrightarrow M$ such
that $F(u)=0$. If every object of $C$ has an injective effacement, see 1.10,
Remark 1, this is equivalent to requiring $F(u)=0$ for every injective
effacement $u$. If every object $A\in C$ has a monomorphism into an injective
object $M$, see 1.10.1, it is equivalent to requiring $F(M)=0$ for every
injective object $M$. Dually, $F$ is called \textit{co-effaceable} if for
every $A\in C$ there exists an epimorphism $u:P\longrightarrow A$ such that
$F(u)=0$.

\LayerTarget{tohoku:prop:2.2.1}
\noindent\textsc{Proposition 2.2.1.} --- \textit{Let $C$ and $C'$ be abelian
categories and let $T=(T^i)_{0\leq i\leq a}$ be an exact
$\partial$-functor, covariant or contravariant, from $C$ to $C'$, where
$a>0$. If $T$ is covariant and $T^i$ is effaceable for $i>0$, then $T$ is a
universal $\partial$-functor. The converse holds if every object $A\in C$ has
an injective effacement; see 1.10. If $T$ is contravariant, the same statement
holds with \textit{co-effaceable} and \textit{projective effacement} in place
of \textit{effaceable} and \textit{injective effacement}, respectively.}

\LayerTarget{tohoku:prop:2.2.1:proof:p1}
For the direct assertion it suffices, for example, to prove the following.
Suppose $(T'^0,T'^1)$ is a $\partial$-functor defined in degrees 0 and 1 and
$f^0:T^0\longrightarrow T'^0$ is a natural transformation. Then there is a
unique morphism $(f^0,f^1)$ from $(T^0,T^1)$ to $(T'^0,T'^1)$ whose degree-zero
component is $f^0$. We assume $T$ covariant to fix ideas. Let $A\in C$ and
choose an exact sequence
$0\longrightarrow A\longrightarrow M\longrightarrow A'\longrightarrow 0$
such that the morphism $T^1(A)\longrightarrow T^1(M)$ is zero. If $f^1$ has
been constructed, there is a commutative diagram

\LayerTarget{tohoku:diagram:prop-2.2.1}
\[
\begin{array}{ccccccc}
T^0(M)&\longrightarrow&T^0(A')&\longrightarrow&T^1(A)&\longrightarrow&T^1(M)\\
\big\downarrow&&\big\downarrow&&\big\downarrow&&\\
T'^0(M)&\longrightarrow&T'^0(A')&\longrightarrow&T'^1(A)&&
\end{array}
\]

\LayerTarget{tohoku:prop:2.2.1:proof:p2}
The first row is exact, so the morphism
$T^0(A')\longrightarrow T^1(A)$ is surjective. Consequently
$f^1(A):T^1(A)\longrightarrow T'^1(A)$ is completely determined on passing
to the quotient by
$f^0(A'):T^0(A')\longrightarrow T'^0(A')$. This proves uniqueness of
$f^1(A)$. Conversely, omit $f^1(A)$ from the preceding diagram. Since the
composite of the two morphisms in the second row is zero, the requirement
that the diagram commute uniquely defines a morphism
$T^1(A)\longrightarrow T'^1(A)$. Standard arguments show that the morphism so
defined is independent of the chosen exact sequence
$0\longrightarrow A\longrightarrow M\longrightarrow A'\longrightarrow 0$,
is natural, and commutes with $\partial$. This proves the first part of the
proposition. The second part follows from the next existence theorem.

\LayerTarget{tohoku:thm:2.2.2:p1}
\noindent\textsc{Theorem 2.2.2.} --- \textit{Let $C$ be an abelian category in
which every object $A\in C$ has an injective effacement; see 1.10. Then, for
every covariant additive functor $F$ on $C$, the satellites $S^iF$ for
$i\geq0$ exist and are effaceable functors for $i>0$. In order that the
universal $\partial$-functor $(S^iF)_{i\geq0}$ be exact, it}

\TohokuContinuation{142}
